\subsection{Quantum couplings of the surviving vectors}
\label{subsec:x9-neutral-quantum}

The smooth-side mirror determines protected vector-multiplet data beyond
the cubic and second-Chern pairings. We compute its genus-zero
Gopakumar--Vafa invariants and compare its electric and neutral magnetic
periods with the continued periods on $\Sigma$. Conifold surgery and
large-volume asymptotics identify these seven periods to all orders on
the marked branch. The finite coefficient calculation provides an
independent check of this identification.

\paragraph{The smooth-side chamber.}
Use the primitive surviving divisor basis
$K_1=H_2$, $K_2=H_3$, $K_3=H_4$, with curve degrees
$\mathbf d=(\ell,m,n)$ defined by their three pairings. Rewriting
\eqref{eq:x9-mutated-relations} in the coordinates $(v,a,b)$ gives
\begin{equation}
 Q^{\rm sm}=\begin{pmatrix}
 1&1&-2&0&0&-1&1\\
 0&0&-1&2&0&1&0\\
 0&0&1&-1&1&0&0
 \end{pmatrix},\qquad Q^{\rm sm}_*=(0,2,1)^T.
 \label{eq:x9-neutral-smooth-glsm}
\end{equation}
Its columns are ordered by $r_0,\ldots,r_6$ and give the toric
divisor classes in the $K$ basis. The only nonsimplicial maximal
cone is $01256$, with circuit $r_0+r_1+r_6=2r_2+r_5$.
Its two simplicial subdivisions replace it, respectively, by
\begin{equation}
 \{1256,0256,0125\},\qquad \{0156,0126\}.
 \label{eq:x9-neutral-ambient-subdivisions}
\end{equation}
These modifications occur over the ambient fixed point avoided by
the general anticanonical hypersurface. They therefore leave
$X_{9,t}$ unchanged. Their ambient Mori cones, in degree coordinates,
are
\begin{equation}
 \begin{aligned}
 M_+&=\operatorname{cone}\{(1,0,0),(0,1,0),(0,0,1)\},\\
 M_-&=\operatorname{cone}\{(-1,0,0),(1,1,0),(0,0,1)\},\\
 \overline{\operatorname{NE}}(X_{9,t})\ &\subseteq\ M_+\cap M_-
 =\operatorname{cone}\{(1,1,0),(0,1,0),(0,0,1)\}.
 \end{aligned}
 \label{eq:x9-neutral-common-mori}
\end{equation}
No equality with the actual Calabi--Yau Mori cone is required.
In either subdivision the ten cubic entries, ordered lexicographically
by triples of $K$ indices, and the second-Chern row are
\begin{equation}
 (0,2,3,6,11,17,14,27,49,83),\qquad (36,80,146).
 \label{eq:x9-neutral-smooth-cubic}
\end{equation}
They coincide with the resolved intersections of $H_2,H_3,H_4$.
This fixes the smooth-side large-volume chamber used for the
calculation. It does not select a crepant model of the birational
mirror small resolution.

\paragraph{The period comparison.}
For the charge matrix \eqref{eq:x9-neutral-smooth-glsm}, the
Frobenius coefficient is
\begin{equation}
 c_{\rm sm}(\ell,m,n)=
 \frac{\Gamma(1+2m+n)}
 {\Gamma(1+\ell)^3\Gamma(1+n-m-2\ell)
  \Gamma(1+2m-n)\Gamma(1+n)\Gamma(1+m-\ell)}.
 \label{eq:x9-neutral-smooth-gamma}
\end{equation}
Its rational derivatives mean derivatives of
$c_{\rm sm}(\mathbf d+\boldsymbol\rho)/
c_{\rm sm}(\boldsymbol\rho)$ at zero, as in
\eqref{eq:x9-mirror-gamma}. In particular, zeros from reciprocal
Gamma factors must be retained before differentiation. Write
$W_{\rm sm},C^{\rm sm}_i,\widehat C^{\rm sm}_{ij}$ for the
resulting scalar, first- and second-derivative series, and put
\begin{equation}
 D^{\rm sm}_i=\frac12\sum_{j,k}\kappa^{\rm sm}_{ijk}
                    \widehat C^{\rm sm}_{jk},\qquad
 g^{\rm sm}_i=C^{\rm sm}_i/W_{\rm sm}.
 \label{eq:x9-neutral-smooth-jets}
\end{equation}
The large-volume mirror construction then applies with
\eqref{eq:x9-neutral-smooth-cubic}
\cite[Section~4.4]{Demirtas:2024Mirror}.

\begin{proposition}\label{prop:x9-neutral-period-descent}
On the marked sufficiently small, nonzero $(v,a,b)$ branch of $\Sigma$,
the scalar, spectator electric and contracted neutral magnetic series
satisfy the analytic identities
\begin{equation}
 \begin{aligned}
 W_{\rm sm}&=W_0|_\Sigma,\\
 C^{\rm sm}_i&=C_{i+1}|_\Sigma,\qquad
 D^{\rm sm}_i=D_{H_{i+1}}|_\Sigma,
       \qquad i=1,2,3.
 \end{aligned}
 \label{eq:x9-neutral-jet-comparison}
\end{equation}
Consequently $\tau_i=t_{i+1}|_\Sigma$ and
\begin{equation}
 \frac{\partial F^{\rm sm}_{\rm inst}}{\partial\tau_i}
 =\Psi_{H_{i+1}}|_\Sigma.
 \label{eq:x9-neutral-exact-instanton-gradient}
\end{equation}
Here $F^{\rm sm}_{\rm inst}$ is the nonpolynomial genus-zero
prepotential defined in \eqref{eq:x9-neutral-instanton-prepotential}.
The statement concerns the specified continuation, not transport
around arbitrary loops in moduli space.
\end{proposition}

\begin{proof}
Let $\check X$ be a nearby smooth member of the original mirror,
and $\check Y$ its projective small resolution at the four-node
locus. The vanishing-cycle space $V\subset H_3(\check X,\QQ)$
is isotropic of rank three by the integral marking. Conifold
surgery identifies the surviving variation of Hodge structure with
\begin{equation}
 H_3(\check Y,\QQ)\simeq V^\perp/V,
 \qquad \dim(V^\perp/V)=14-2\cdot3=8.
 \label{eq:x9-neutral-period-quotient}
\end{equation}
Periods of classes in $V^\perp$ restrict to the corresponding
periods of $\check Y$; their values on $V$ vanish. This is the
weight-three quotient in the limiting mixed Hodge structure
\cite[Corollary~1.11 and Section~3.3]{Lee:2018Conifold}.
All electric classes and the magnetic classes with divisor part
$H_2,H_3,H_4$ lie in $V^\perp$: the vanishing charges are pure
D2 charges, and these three divisors annihilate them.

Proposition~\ref{prop:x9-mirror-mutation} identifies $\check Y$
birationally with the smooth-side Batyrev mirror and preserves the
residue form. The two smooth models are connected by flops.
Each flop is an isomorphism away from curves; three-cycles can be
moved off those curves, so their period integrals are preserved.
Applied over a local parameter neighborhood, this identifies their
period spaces with the Gauss--Manin connection. Thus the restricted
neutral periods and the smooth-side Frobenius periods belong to
the same eight-dimensional space of local period functions.
The scalar identity fixes their common normalization of the form.

It remains to identify the particular functions within that space.
This step uses unnormalized periods, not the nonlinear regular
corrections $\Psi_P$. Put
$\lambda=(\log v,\log a,\log b)/(2\pi i)$.
For either collection of series $(W,C_i,D_i)$, define
\begin{equation}
 \begin{aligned}
 \Pi^{\rm el}_i&=W\lambda_i+\frac{C_i}{2\pi i},\\
 \Pi^{\rm mag}_i&=\frac W2\sum_{j,k}\kappa^{\rm sm}_{ijk}
                         \lambda_j\lambda_k
   +\frac1{2\pi i}\sum_{j,k}\kappa^{\rm sm}_{ijk}\lambda_j C_k
   +\frac{D_i}{(2\pi i)^2}.
 \end{aligned}
 \label{eq:x9-neutral-linear-periods}
\end{equation}
These are linear combinations of Frobenius periods. Subtracting
the universal $\zeta(2)$ term in the rational second derivative
only adds a constant multiple of $W$. For the resolved periods
restricted to $\Sigma$, the local logarithms and the three local
$C_a$ vanish. Their remaining cubic contractions are exactly
\eqref{eq:x9-neutral-smooth-cubic}, giving
\eqref{eq:x9-neutral-linear-periods} with
$C_i=C_{i+1}|_\Sigma$ and $D_i=D_{H_{i+1}}|_\Sigma$.

The constant terms of both collections are
\begin{equation}
 W(0)=1,\qquad C_i(0)=D_i(0)=0.
 \label{eq:x9-neutral-boundary-constants}
\end{equation}
On the resolved side, every unbounded second-derivative sector
is multiplied by a zero neutral pairing. All remaining local
indices at spectator degree zero are zero, where the rational
Gamma coefficient is identically one. This proves the last
equality without truncating an infinite tail. The previously
proved convergence and the smooth-side Frobenius expansion give
the common polynomial asymptotics
\begin{equation}
 W\sim1,\qquad \Pi^{\rm el}_i\sim\lambda_i,\qquad
 \Pi^{\rm mag}_i\sim\frac12\kappa^{\rm sm}_{ijk}
                                  \lambda_j\lambda_k.
 \label{eq:x9-neutral-logarithmic-asymptotics}
\end{equation}
For example, on $(v,a,b)=(e^{-s w_1},e^{-s w_2},e^{-s w_3})$
with every $w_i>0$, the omitted terms are bounded by
$O(e^{-s\min w_i}(1+s)^2)$. All three weights may vary; this
is a large-volume comparison at the origin, not the circle
path with $v$ fixed.

The smooth-side period space has a Frobenius basis with polynomial
terms of degrees $0,1,1,1,2,2,2,3$
\cite[Equations~(4.15)--(4.19)]{Demirtas:2024Mirror}.
Its three leading quadratics are those in
\eqref{eq:x9-neutral-logarithmic-asymptotics}. They are independent,
as witnessed by
\begin{equation}
 \det\begin{pmatrix}0&2&3\\2&6&11\\3&11&17\end{pmatrix}=10.
 \label{eq:x9-neutral-logarithmic-minor}
\end{equation}
The rows are $(\kappa_{i11},\kappa_{i12},\kappa_{i13})$.
The eighth period has nonzero leading cubic
$\kappa^{\rm sm}_{ijk}\lambda_i\lambda_j\lambda_k/6$,
whose $\lambda_3^3$ coefficient is $83/6$.
Therefore the eight polynomial asymptotics are linearly independent.
No nonzero constant linear combination of these periods can be
exponentially small for every positive weight vector: successively
its cubic, quadratic, linear and constant polynomial coefficients
would have to vanish.

Subtract the smooth-side functions in
\eqref{eq:x9-neutral-linear-periods} from their restricted resolved
counterparts. Each difference is a period function in this common
space and has zero polynomial asymptotic by
\eqref{eq:x9-neutral-boundary-constants}--\eqref{eq:x9-neutral-logarithmic-asymptotics}.
Each difference is consequently zero. The electric identities first
give equality of the $C_i$, and the magnetic identities then give
equality of the $D_i$. The master equation
\eqref{eq:x9-mirror-master} proves
\eqref{eq:x9-neutral-exact-instanton-gradient}.
\end{proof}

An independent coefficient check sums the resolved ten-monomial
series over every local degree for all $35$ triples with
$\ell+m+n\leq4$, and verifies
\eqref{eq:x9-neutral-jet-comparison} in exact arithmetic.
The four unbounded second-derivative sectors of
\eqref{eq:x9-magnetic-tail-lines} are multiplied by
$K_i\cdot C_1$, $K_i\cdot C_2$, $K_i\cdot\iota_*e_1$ or
$K_i\cdot\iota_*e_2$, all zero. Every other nonzero sector has
local indices at most $4(\ell+m+n)$, so each tested neutral
coefficient includes its complete local-degree sum. The finite check does not replace the period-space argument.
The proposition identifies the three instanton gradients, and hence
all genus-zero Yukawa couplings after the already established cubic
match. This argument leaves a moduli-independent prepotential constant.
Section~\ref{subsec:x9-d6-period} fixes it and the remaining D6 period
using the polarization and the reality of the marked continuation,
with the curvature and flux shifts of an integral symplectic basis.
No equality of all uncontracted second Gamma derivatives is asserted.

\paragraph{Protected curve invariants.}
Define $q_i=\exp(2\pi i\tau_i)$ by
$2\pi i\tau_i=\log(v,a,b)_i+g^{\rm sm}_i$.
The nonpolynomial genus-zero prepotential is
\begin{equation}
 F^{\rm sm}_{\rm inst}
 =-\frac1{(2\pi i)^3}\sum_{\mathbf d\ne0}
       n^{0,\rm sm}_{\mathbf d}\operatorname{Li}_3(q^{\mathbf d}).
 \label{eq:x9-neutral-instanton-prepotential}
\end{equation}
Applying the three magnetic equations
\eqref{eq:x9-mirror-master} with $Q^{\rm sm}$ determines
$n^{0,\rm sm}_{\mathbf d}$ for all $164$ nonzero nonnegative
degree vectors with $\ell+m+n\leq8$. All three equations agree,
every invariant is integral, and $31$ are nonzero. The complete
nonzero list through degree three is
\begin{equation}
 \begin{array}{c|rrrrr}
 \mathbf d&(0,0,1)&(1,1,0)&(0,1,1)&(1,1,1)&(0,1,2)\\\hline
 n^{0,\rm sm}_{\mathbf d}&10&3&140&-20&140
 \end{array}.
 \label{eq:x9-neutral-gv-low}
\end{equation}
The degree is taken in the primitive $K$ basis, not in
the nonnegative-relation coordinates $\xi_i$.

The plane divisor is the $r_2$ column,
$\nu=-2K_1-K_2+K_3$. Its square pairs as $(-3,-3,0)$,
so a plane line has primitive degree $(1,1,0)$. Along this ray
the compact calculation gives
\begin{equation}
 n^{0,\rm sm}_{(d,d,0)}=3,-6,27,-192,
 \qquad d=1,2,3,4.
 \label{eq:x9-neutral-plane-gv}
\end{equation}
An independent one-parameter calculation with local
$\mathcal O_{\mathbb P^2}(-3)$ charges $(1,1,1,-3)$ reproduces
both these invariants and the corresponding coefficients of
$g^{\rm sm}_1+g^{\rm sm}_2$. These are genus-zero protected
M2 indices on the smooth side \cite{Gopakumar:1998BPS}; they are
not unweighted state counts, magnetic probe indices, or the
spectrum of the cooperative condensate at the transition.
